%% ma: L11-B04-0018
%% lop: 11
%% chuong: 1
%% bai: 4
%% dang: 11.4.1
%% loai: TL
%% muc: TH
%% dap_an: "-"
%% nguon: "de-on-tap-giua-ki-1-toan-11-nam-2026-2027-truong-thpt-ngo-quyen-da-nang.pdf (M09, Đề ôn tập giữa kì 1 Toán 11 THPT Ngô Quyền, Đà Nẵng), Mã đề 0111, Phần IV Câu 1"
%% kien_thuc: "Bài 3 (hàm số lượng giác), Bài 4 (phương trình lượng giác cơ bản)"
%% trang_thai: chinh-thuc
\begin{ex}
Mực nước $h(t)$ (tính bằng mét) tại một cảng biển ở thời điểm $t$ giờ sau nửa đêm $(0\le t\le24)$ được mô hình hóa bởi công thức $h(t)=3+2\cos\left(\dfrac{\pi t}{6}\right)$.

a) Tính độ sâu lớn nhất của nước tại cảng biển và xác định thời điểm đầu tiên trong ngày mà mực nước đạt độ sâu lớn nhất.

b) Trong khoảng thời gian từ $0$ giờ đến $12$ giờ trưa, xác định các thời điểm mà mực nước tại cảng có độ sâu đúng bằng $4$ m.
\loigiai{a) $h(t)\le3+2=5$, dấu bằng khi $\cos\dfrac{\pi t}6=1\Leftrightarrow t=12k$. Độ sâu lớn nhất là $5$ m, đạt lần đầu tiên lúc $t=0$ ($0$ giờ, nửa đêm).

b) $h(t)=4\Leftrightarrow\cos\dfrac{\pi t}6=\dfrac12\Leftrightarrow\dfrac{\pi t}6=\pm\dfrac\pi3+k2\pi\Leftrightarrow t=\pm2+12k$. Với $0\le t\le12$: $t=2$ hoặc $t=10$, tức là lúc $2$ giờ sáng và $10$ giờ sáng.}
\end{ex}
